%% GENERATED by tools/gen_latex.py -- edit content/viny-llm.py instead.
%% Build:  latexmk -pdf viny-llm.tex     (or: pdflatex viny-llm.tex, twice)
\documentclass[aspectratio=169,11pt,t]{beamer}
\usepackage{itbpro}

\coursetitle{Large Language Models}
\coursesubtitle{Fine-tuning, retrieval, re-ranking, and guardrails --- the practitioner's layer on top of chapters 15 and 16.}
\coursekicker{Module}
\coursesource{Module material for Designing and Building AI Products and Services}
\courseduration{3 hours (2 sessions)}
\coursepresenter{Viny Christanti Mawardi, S.Kom., M.Kom.}{Teaching Assistant}
\coursebrand{AI for Professional \textperiodcentered\ ITB}

\title{Large Language Models}
\author{Viny Christanti Mawardi, S.Kom., M.Kom.}
\date{}

\begin{document}
\covertitle

\framekicker{Learning outcomes}
\begin{frame}[shrink=25]{Session Objectives}
\leadpar{By the end of this session, participants can:}
\begin{isteps}
\item \textit{Objectives pending --- to be written by the module owner.}
\item Placeholder: explain the \textbf{pretraining / adaptation / retrieval} division of labour.
\item Placeholder: choose between \textbf{full fine-tuning, LoRA, and prompting} for a given problem.
\item Placeholder: build a \textbf{retrieval pipeline} with a re-ranker and measure it.
\item Placeholder: specify \textbf{guardrails} for a production deployment.
\end{isteps}

\vskip 6pt
\vskip 4pt
\reslink{Site}{Course home}{https://hendrikarisma.my.id/rs/ai-products-course/index.html}
\vskip 2pt
\end{frame}

\framekicker{About this deck}
\begin{frame}[shrink=25]{This is a template, not the module}
\begin{iband}[amber]
\textbf{This deck is scaffolding.} The structure, theme, and rendering are fixed; the material is authored separately by the module owner and may be delivered as PDF or PPTX instead.
\end{iband}
{\fontsize{9.4}{12.6}\selectfont\color{fg2}It exists so the course site has a real deck to link to from day one, and so that --- if the module is authored here --- every block kind it will plausibly need is already demonstrated and verified in both the LaTeX and the web output.\par}\vskip 4pt
{\fontsize{9.4}{12.6}\selectfont\color{fg2}Slides carrying no final content are marked with an amber band, so that an unfinished deck \textbf{cannot be mistaken for a finished one} on screen.\par}\vskip 4pt
\end{frame}
\note{Delete this slide once the module has real content.}

\framekicker{Where this module sits}
\begin{frame}[shrink=25]{On top of chapters 15 and 16}
{\fontsize{9.4}{12.6}\selectfont\color{fg2}Chapter 15 built the Transformer and fine-tuned an encoder. Chapter 16 built a mini-GPT, sampled from it, and instruction-tuned a pretrained model under LoRA.\par}\vskip 4pt
{\fontsize{9.4}{12.6}\selectfont\color{fg2}This module takes all of that as given and covers the \textbf{practitioner's layer}: what you do when the model is someone else's, the data is yours, and the answer has to be right.\par}\vskip 4pt
\begin{center}
\includegraphics[width=0.92\linewidth,height=0.60\textheight,keepaspectratio]{figs/viny-arc.tex.pdf}
\end{center}
\vskip 3pt{\color{ink3}\fontsize{7.4}{9.4}\selectfont Placeholder outline --- replace with the module's own structure.\par}
\end{frame}

\coursesection{01}{Pretraining and adaptation}{Section lead pending.}

\framekicker{Section 01}
\begin{frame}[shrink=25]{Adaptation options, side by side}
{\fontsize{9.4}{12.6}\selectfont\color{fg2}A comparison table is usually the right opening for this material. The row structure below is a suggestion; the entries are placeholders.\par}\vskip 4pt
\begin{itable}{>{\raggedright\arraybackslash}p{0.2444\textwidth}>{\raggedright\arraybackslash}p{0.3572\textwidth}>{\raggedright\arraybackslash}p{0.3384\textwidth}}
\thead{Approach} & \thead{When it fits} & \thead{What it costs} \\ \midrule
\textbf{Prompting} & \textit{pending} & \textit{pending} \\
\textbf{Few-shot prompting} & \textit{pending} & \textit{pending} \\
\textbf{LoRA fine-tuning} & \textit{pending} & \textit{pending} \\
\textbf{Full fine-tuning} & \textit{pending} & \textit{pending} \\
\textbf{Continued pretraining} & \textit{pending} & \textit{pending} \\
\end{itable}
\begin{iband}[amber]
\textbf{Content pending.} This slide is scaffolding --- the structure is fixed, the material is authored separately by the module owner.
\end{iband}
\end{frame}

\framekicker{Section 01 $\cdot$ block demo}
\begin{frame}[shrink=25]{Code and output render like this}
{\fontsize{9.4}{12.6}\selectfont\color{fg2}A code block always needs prose before and after it --- the build lints for this, and the check applies to module decks too.\par}\vskip 4pt
\icode{python}{placeholder}{listings/viny-llm-01.py}
\ioutput{listings/viny-llm-02.txt}
{\fontsize{9.4}{12.6}\selectfont\color{fg2}Program output uses the \inlinecode{out} block, which renders in a distinct panel so a result is never mistaken for source.\par}\vskip 4pt
\end{frame}

\coursesection{02}{Retrieval and re-ranking}{Section lead pending.}

\framekicker{Section 02}
\begin{frame}[shrink=25]{Where the re-ranker sits}
{\fontsize{9.4}{12.6}\selectfont\color{fg2}Chapter 16 introduced RAG as \textit{retrieve, then put it in the prompt}. The practitioner's version has one more stage, and it is usually where the quality comes from.\par}\vskip 4pt
\begin{center}
\includegraphics[width=0.92\linewidth,height=0.60\textheight,keepaspectratio]{figs/viny-rag.tex.pdf}
\end{center}
\vskip 3pt{\color{ink3}\fontsize{7.4}{9.4}\selectfont The retriever optimises recall cheaply; the re-ranker optimises precision expensively, over far fewer candidates.\par}
\begin{iband}[amber]
\textbf{Content pending.} This slide is scaffolding --- the structure is fixed, the material is authored separately by the module owner.
\end{iband}
\end{frame}

\framekicker{Section 02 $\cdot$ block demo}
\begin{frame}[shrink=25]{Cards, stats, and quotes render like this}
\begin{minipage}[t]{0.3253\textwidth}
\istat{---}{recall@k, pending}
\end{minipage}%
\hfill
\begin{minipage}[t]{0.3253\textwidth}
\istat{---}{nDCG, pending}
\end{minipage}%
\hfill
\begin{minipage}[t]{0.3253\textwidth}
\istat{---}{latency, pending}
\end{minipage}%
\begin{minipage}[t]{0.3253\textwidth}
\begin{icardaccent}
\cardh{Dense retrieval}
\textit{pending}
\end{icardaccent}
\end{minipage}%
\hfill
\begin{minipage}[t]{0.3253\textwidth}
\begin{icardaccent}
\cardh{Sparse retrieval}
\textit{pending}
\end{icardaccent}
\end{minipage}%
\hfill
\begin{minipage}[t]{0.3253\textwidth}
\begin{icardaccent}
\cardh{Hybrid}
\textit{pending}
\end{icardaccent}
\end{minipage}%
\begin{iquote}
\quotetext{A pull quote renders like this --- useful for a definition or a finding worth stopping on.}\quotecite{Attribution line}
\end{iquote}
\end{frame}

\coursesection{03}{Guardrails and evaluation}{Section lead pending.}

\framekicker{Section 03}
\begin{frame}[shrink=25]{What has to be checked, and where}
\begin{isteps}
\item \textit{Step one --- pending.}
\item \textit{Step two --- pending.}
\item \textit{Step three --- pending.}
\end{isteps}
\begin{iband}[amber]
\textbf{Content pending.} This slide is scaffolding --- the structure is fixed, the material is authored separately by the module owner.
\end{iband}
{\fontsize{9.4}{12.6}\selectfont\color{fg2}Numbered steps use the \inlinecode{steps} block; unordered points use \inlinecode{bullets}. Both wrap correctly in the LaTeX output at any length.\par}\vskip 4pt
\end{frame}

\framekicker{Summary}
\begin{frame}[shrink=25]{Summary slide --- pending}
\begin{isteps}
\item \textit{Takeaway one --- pending.}
\item \textit{Takeaway two --- pending.}
\item \textit{Takeaway three --- pending.}
\end{isteps}
\vskip 4pt
\reslink{BACK}{Chapter 16 --- Text generation}{https://hendrikarisma.my.id/rs/ai-products-course/ch16/index.html}
\reslink{NEXT}{Agentic AI}{https://hendrikarisma.my.id/rs/ai-products-course/hendri-agentic/index.html}
\vskip 2pt
\end{frame}

\end{document}
